\section{Threats to Validity}
\label{sec:threats}

For each threat we state the mitigation we actually performed and the risk that remains. Threats that we could only measure or disclose are collected at the end of the section.

\subsection{Internal Validity}

\textbf{Reverse-engineered stories.} The user stories were written by GPT-5.4-mini from the reference files, so the input can carry wording, structure and scenario order from the very file it is compared with, and the model that wrote the stories is also the model under test. \emph{Mitigation:} we measured story--reference trigram overlap (median 4.9\%) and its correlation with similarity (Section~\ref{sec:res-b1}), and we checked the 84 Fineract stories that were written from only the first 15,000 characters of a longer reference; they score no lower than the other Fineract files (Section~\ref{sec:res-h1}). \emph{Residual risk:} overlap correlates with similarity in two of the three projects, so part of the measured similarity likely comes from this construction; structural leakage such as the order and number of criteria is not measured; and no independently written backlog was available, so the results may overstate performance on natural requirements.

\textbf{Exemplars and pilot.} \emph{Mitigation:} the three original exemplars were part of the dataset and are excluded from the primary analysis; both denominators are reported. \emph{Residual risk:} the 30 pilot stories used when the reference point was set are part of the evaluated set.

\textbf{Prompt ablation.} The original ablation ran three months after C0. \emph{Mitigation:} the same-time rerun repeats C0, C1 and C2 with interleaved requests, and its C0 has the same median as the original C0. \emph{Residual risk:} C2 uses one random draw of exemplars, and the rerun cannot show what changed in the hosted service between the two dates.

\textbf{Hosted model.} The serving infrastructure is outside our control. \emph{Mitigation:} we pinned the snapshot, used temperature 0, logged each call, recorded the returned model name for C1, C2 and the rerun, and measured run-to-run variation on 60 files (Section~\ref{sec:res-rep}). \emph{Residual risk:} nondeterminism remains, the original C0 log records only the requested name, and the API returned no system fingerprint, so changes of the serving configuration cannot be detected.

\textbf{Output cap.} The proposal planned a cap of 1,500 output tokens that the original runs did not apply. \emph{Mitigation:} the rerun included this capped condition (C0cap). On the 258 non-exemplar files, 81 outputs stopped at the cap, and 239 files passed the parser (92.6\%, 95\% CI [88.7\%, 95.5\%]) against 258 for the rerun C0 (exact McNemar $p < 0.001$); all 19 failures are capped outputs. The median cosine was 0.854, with a median paired difference of $+0.0005$ (Wilcoxon $p = 0.108$). \emph{Residual risk:} the RQ2 result applies to uncapped generation; under the cap, long files are cut and parser compliance is lower. These two tests are descriptive and uncorrected.

\subsection{Construct Validity}

\textbf{Similarity metric.} Skeleton cosine is an imperfect proxy for business correctness. \emph{Mitigation:} we checked it against defect labels, seeded defects, reference points from 0.80 to 0.90 and per project. \emph{Residual risk:} it is blind to numbers and quoted values by construction, ignores data tables and Examples, and for about two thirds of the files it sees only the first 256 tokens of a skeleton. Within the audited high-similarity files it did not separate defective from non-defective outputs significantly, scenario-level seeded defects barely changed it, and a valid alternative scenario can score low because it differs from the single project-maintained reference.

\textbf{Defect labels.} The 60 audit labels were drafted by an AI assistant and retained unchanged by a reviewing author, whose review was recorded as one confirmation rather than per-file decisions. \emph{Mitigation:} an independent human rater labelled 15 files (independence rests on the rater's own declaration), and we report $\kappa$ with 95\% intervals. Because $\kappa$ was below our protocol's threshold, the reviewing author and the independent rater then held the required consensus round on these 15 files (Section~\ref{sec:res-defects}). \emph{Residual risk:} $\kappa = 0.571$ (business logic) and $0.609$ (temporal/causal) are below the 0.7 our protocol required, and with 15 files they are imprecise (the intervals reach down to 0.05 and 0.24). The protocol also required relabelling, but no rater who had not seen the labels was available for the other 45 files, so we keep their reviewed labels, report $\kappa$ as measured before consensus and treat the audit as an exploratory observation rather than as validation. The consensus resolved every disagreement in favour of the independent rater and added defects in three Diaspora files, so the reviewed labels of the other 45 files may undercount defects, in particular in the social-network domain. The two label facets are not independent measurements, and the audit covers only files at or above 0.85.

\textbf{Perturbation test.} \emph{Mitigation:} pairs in which a text does not parse are excluded, and results are reported by operator (Appendix~\ref{app:perturb}). \emph{Residual risk:} the test mutates single extracted scenarios, omitted the Background of 74 of 120 source scenarios, and draws several mutants from the same scenario, so its pairs are not independent and do not estimate whole-file sensitivity.

\textbf{Parser compliance.} \emph{Mitigation:} the submitted version's dry-run check, which counted a configuration error as a PASS, was replaced by behave's parser, and all parser results were recomputed. \emph{Residual risk:} a parser pass shows only that a file is valid Gherkin, not successful execution against the application or correct step implementations.

\subsection{External and Conclusion Validity}

\textbf{Projects.} \emph{Mitigation:} all main results are reported per project. \emph{Residual risk:} the dataset spans three mature, English-language public projects and is not representative of BDD practice; results may differ for short tickets, mixed-language requirements, private conventions or projects with unstable step definitions. The projects contribute unequal numbers of files (99, 99 and 60 after exclusion) and differ significantly in similarity, so pooled results do not generalise across domains.

\textbf{Thresholds and tests.} \emph{Mitigation:} we report how the conclusions change between reference points of 0.80 and 0.90, and the ablation and rerun plans were fixed with their hashes before the runs. \emph{Residual risk:} neither the 0.85 reference point nor the 80\% threshold is an industry standard, the reference point was not calibrated, and the RQ1 test result changes between 0.84 and 0.85. All analyses added in revision were planned after the reviews and are exploratory; their $p$-values are not corrected across analyses.

\textbf{Uncontrolled limitations.} The measures above do not control every threat. We had no independently written requirements, could not detect changes in the hosted service, did not execute the generated files against step definitions, and could not relabel the 45 audit files outside the consensus round. The conclusions therefore hold only for reverse-engineered stories from these three projects, in the time window of the runs, and with the similarity metric read as a proxy.
